%% ma: L10-B01-0009
%% lop: 10
%% chuong: 1
%% bai: 1
%% dang: 10.1.4
%% loai: TN4
%% muc: TH
%% dap_an: "D"
%% nguon: "Ảnh giáo viên gửi 10/10/2026 (ThS. Nguyễn Hoàng Việt, Chương 1 Mệnh đề), A-Đề 1, Phần 1 Câu 10"
%% kien_thuc: "Bài 1 (mệnh đề chứa kí hiệu ∀, ∃)"
%% trang_thai: cho-duyet
%% ly_do: "Dạng toán mới đề xuất 10.1.4: mệnh đề chứa kí hiệu ∀, ∃: xét tính đúng sai"
\begin{ex}
Cho $n$ là số tự nhiên, mệnh đề nào sau đây đúng?
\choice{$\forall n,\ n(n+1)$ là số chính phương}{$\forall n,\ n(n+1)$ là số lẻ}{$\exists n,\ n(n+1)(n+2)$ là số lẻ}{\True $\forall n,\ n(n+1)(n+2)$ chia hết cho $6$}
\loigiai{Trong ba số tự nhiên liên tiếp luôn có một số chẵn và một số chia hết cho $3$, nên tích chia hết cho $2$ và $3$, tức chia hết cho $6$: D đúng. A sai ($n=1$: $2$ không chính phương); B sai vì $n(n+1)$ luôn chẵn; C sai vì $n(n+1)(n+2)$ luôn chẵn.}
\end{ex}
